\documentclass[11pt]{article}

\usepackage[utf8]{inputenc}
\usepackage[T1]{fontenc}
\usepackage{mathptmx}
\usepackage[margin=1in]{geometry}
\usepackage{booktabs}
\usepackage{array}
\usepackage[hidelinks]{hyperref}
\usepackage{float}
\usepackage{titlesec}

\setlength{\parskip}{0.7em}
\setlength{\parindent}{0pt}
\titlespacing*{\section}{0pt}{1.2ex plus 0.5ex minus 0.2ex}{0.6ex plus 0.2ex}
\titlespacing*{\subsection}{0pt}{1.0ex plus 0.5ex minus 0.2ex}{0.4ex plus 0.2ex}

\title{Study B: Does a DP-Protected Sex Label Preserve the Subgroup AUC Gap? \\ \large Draft --- Findings and Numbers}
\author{Andrew}
\date{August 19, 2026}

\begin{document}
\maketitle

\begin{abstract}
\noindent This is an internal, numbers-focused draft covering all Study B results to date. Study B asks whether a differentially private release of patient sex preserves or washes out the subgroup pneumothorax AUC gap established in Study A, on the frozen 90/10 predictions file, across a fixed epsilon sweep. The mechanism used to answer this question was reworked once during development: an initial aggregate-count-based design was discarded after a technical review found it did not actually produce privacy-protected data (Section~\ref{sec:mechanism-history}). The corrected mechanism --- per-record randomized response --- was itself only single-draw-evaluated at first; a subsequent 30-seed replication (Section~\ref{sec:replication}) found that the single canonical draw's headline claims about direction reliability and about epsilon=4--5 being stable were both artifacts of that one draw, not properties of those epsilon values. The replication-corrected finding is that the gap is fully reliable only at epsilon $\geq6$; below that, survival is probabilistic and, at epsilon $\leq1$, direction itself is unreliable roughly 10--30\% of the time. No narrative analysis beyond what is needed to interpret the numbers is included here; this is a results record, not a paper section.
\end{abstract}

\section{Methods (brief)}

\textbf{Input.} Study A's frozen \texttt{predictions\_90\_10.csv} and \texttt{patient\_split.csv} (test split only): 4{,}620 test patients, 2{,}470 male / 2{,}150 female, patient-level \texttt{true\_label}/\texttt{predicted\_score} aggregated by the same max/mean rule Study A uses. True gap (male AUC $-$ female AUC), recomputed directly from the frozen CSV as a cross-check against Study A's own reported value: 0.067036, matching Study A's canonical 0.0670.

\textbf{Mechanism.} \texttt{privatize\_categorical\_label} (\texttt{src/dp\_mechanisms.py}): diffprivlib's \texttt{Binary} mechanism, classic randomized response, applied independently to every test patient's true sex label. Each label is kept with probability $e^\epsilon/(1+e^\epsilon)$ and flipped otherwise --- a single, well-defined per-record epsilon-DP query, verified empirically to match that formula exactly.

\textbf{Epsilon sweep.} $\{0.1, 0.5, 1, 2, 3, 4, 5, 6, 8, 10\}$ --- re-derived specifically for this mechanism's noise profile (a diagnostic found the transition zone at epsilon 3--4, comfortably inside the original six commonly-cited values, unlike the discarded mechanism, which never washed out the gap anywhere in that range).

\textbf{Test oracle.} At each epsilon, recompute the subgroup AUC gap using the per-record DP-protected sex labels directly (\texttt{true\_label}/\texttt{predicted\_score} unchanged). ``Survives'' = same direction and magnitude within 15\% of the true gap. ``Washed out'' = wrong direction or $>$15\% off. The 15\% threshold and the epsilon sweep are both fixed by project convention, not derived from this data.

\section{Results}

\subsection{Mechanism history: a discarded first version}
\label{sec:mechanism-history}

The first implementation released DP-noised aggregate male/female counts once per epsilon, then reconstructed a per-patient group assignment by randomly moving just enough patients between groups --- starting from their \emph{true} labels --- to match the noised aggregate target. A technical review before pushing found two disqualifying problems, not minor caveats:

\begin{enumerate}
\item \textbf{Construct validity.} At epsilon=10, 0 of 4{,}620 patients were ever reassigned; even at epsilon=0.1, only $\sim$7 were. The artifact scored for AUC was true labels with minimal random perturbation, not privacy-protected data --- the experiment could not have shown a different qualitative result across almost the entire sweep regardless of whether genuinely privacy-protected per-record data would preserve the audit.
\item \textbf{Privacy accounting.} The aggregate mechanism's ``full epsilon per category via parallel composition'' claim holds under add/remove adjacency, but the test cohort's membership is already public (\texttt{patient\_split.csv} is committed to git) --- the real threat model is attribute (substitute-one) privacy, under which that composition claim does not hold (one attribute flip changes two bin counts simultaneously, breaking parallel composition's precondition).
\end{enumerate}

All commits implementing this version were discarded (the branch had never been pushed). The replacement mechanism (randomized response, above) resolves both problems: every output label is a genuine per-record draw, and there is no aggregate count or cross-bin composition question at all.

\subsection{Canonical sweep (single draw, seed 42)}

\begin{table}[H]
\centering
\begin{tabular}{rrrrl}
\toprule
Epsilon & DP gap & Pct.\ diff.\ from true gap & Direction match & Survived \\
\midrule
0.1 & 0.0355 & 47.1\% & True & False \\
0.5 & 0.0572 & 14.7\% & True & True \\
1.0 & 0.0466 & 30.4\% & True & False \\
2.0 & 0.0936 & 39.7\% & True & False \\
3.0 & 0.0596 & 11.1\% & True & True \\
4.0 & 0.0661 & 1.3\% & True & True \\
5.0 & 0.0672 & 0.3\% & True & True \\
6.0 & 0.0660 & 1.5\% & True & True \\
8.0 & 0.0663 & 1.1\% & True & True \\
10.0 & 0.0670 & 0.0\% & True & True \\
\bottomrule
\end{tabular}
\caption{Canonical single-draw sweep (\texttt{BASE\_SEED=42}). True gap = 0.0670 throughout. Crossover epsilon (smallest value at/above which every larger epsilon also survives): 3.0.}
\label{tab:canonical-sweep}
\end{table}

Read on its own, this table appears to show a clean story: direction always correct, and a fairly sharp, stable transition around epsilon=3--4. Section~\ref{sec:replication} shows this reading does not survive replication.

\subsection{30-seed replication}
\label{sec:replication}

30 independent draws per epsilon, a seed pool separate from the canonical sweep, $\sim$80s total runtime.

\begin{table}[H]
\centering
\begin{tabular}{rrrrr}
\toprule
Epsilon & Mean DP gap (SD) & Mean pct.\ diff. & Direction agreement & Survival rate \\
\midrule
0.1 & 0.0119 (0.0269) & 82.7\% & 70.0\% & 3.3\% \\
0.5 & 0.0169 (0.0240) & 75.0\% & 73.3\% & 10.0\% \\
1.0 & 0.0232 (0.0236) & 65.8\% & 90.0\% & 6.7\% \\
2.0 & 0.0504 (0.0164) & 26.0\% & 96.7\% & 36.7\% \\
3.0 & 0.0626 (0.0108) & 14.1\% & 100\% & 50.0\% \\
4.0 & 0.0652 (0.0073) & 8.3\% & 100\% & 90.0\% \\
5.0 & 0.0651 (0.0047) & 4.1\% & 100\% & 90.0\% \\
6.0 & 0.0668 (0.0016) & 1.7\% & 100\% & 100\% \\
8.0 & 0.0671 (0.0004) & 0.3\% & 100\% & 100\% \\
10.0 & 0.0671 (0.0001) & 0.04\% & 100\% & 100\% \\
\bottomrule
\end{tabular}
\caption{30-replicate summary per epsilon, separate seed pool (\texttt{REPLICATION\_BASE\_SEED=2000}) from the canonical sweep. Replication-based crossover epsilon (smallest value at/above which survival rate $\geq$50\% holds for every larger epsilon too): 3.0, numerically the same as the canonical draw's, but see Section~\ref{sec:discussion} for why the per-epsilon story is materially different.}
\label{tab:replication}
\end{table}

\section{Discussion}
\label{sec:discussion}

\textbf{The replication changed the headline, not just the error bars.} Two claims that read cleanly off the single canonical draw (Table~\ref{tab:canonical-sweep}) do not hold up:

\begin{itemize}
\item \emph{Direction is unreliable at low epsilon, not just magnitude.} The canonical draw showed \texttt{direction\_match=True} at every epsilon, including 0.1. Across 30 replicates, direction agreement is only 70\% at epsilon=0.1, 73\% at 0.5, 66--90\% through epsilon=1--2 (Table~\ref{tab:replication}). The single draw getting direction right at low epsilon was luck, not a property of that epsilon.
\item \emph{Epsilon=4 and 5 are not the stable points the canonical draw suggested.} The canonical draw reported \texttt{survived=True} with pct.\ diff.\ of 1.3\% and 0.3\% at epsilon=4 and 5 respectively --- numbers that look like a settled result. The replication shows only 90\% survival at each: 1 in 10 draws at those exact epsilon values does not survive. Epsilon=3 is a near-exact coin flip (50.0\% survival). Only epsilon $\geq6$ reaches 100\% survival across all 30 replicates.
\end{itemize}

\textbf{Numeric agreement on the crossover value does not mean the underlying claims agree.} Both the canonical single draw and the replication-based definition land on crossover epsilon = 3.0. That agreement is coincidental at the level of a single summary statistic; it does not rescue the canonical draw's specific claims about epsilon=4 and 5, which the replication shows were overstated.

\textbf{What is robust.} The gap's \emph{direction} is essentially always preserved once epsilon $\geq2$ (96.7--100\% agreement), and magnitude is fully reliable (100\% survival, SD $\leq$0.002) at epsilon $\geq6$. This is a genuine, defensible finding: at privacy budgets on the looser end of what's commonly cited in practice, DP-protected sex labels reliably preserve this subgroup bias-audit conclusion on a cohort this size.

\textbf{What is not robust.} Anything in the epsilon=0.5--5 range. Survival rate climbs roughly monotonically with epsilon in that band (10\%, 6.7\%, 36.7\%, 50\%, 90\%, 90\% at 0.5/1/2/3/4/5) but is nowhere near deterministic until epsilon=6. A policy claim of the form ``DP survives at epsilon=$X$'' for any $X$ in this range should cite the replication's survival rate at that specific epsilon, not a single-draw boolean.

\section{Limitations}
\begin{itemize}
\item Scope is the single 90/10 arm and Pneumothorax only, per Study B's fixed design (it does not sweep imbalance ratios or read Study A's other prediction files). Nothing here establishes whether this mechanism's reliability threshold would look different against a smaller true gap (e.g.\ Study A's 70/30 or 50/50 arms, whose gaps are 0.017--0.04 --- close enough to the evaluation noise floor that a 15\%-of-true-gap tolerance would be far tighter in absolute AUC terms).
\item The 15\% survival threshold and the epsilon sweep values are both fixed by project convention (CLAUDE.md), not derived from or validated against this data; no sensitivity analysis on the 15\% figure has been run.
\item No analytical or closed-form relationship between the randomized-response flip probability and expected AUC-gap attenuation has been derived --- every number in Table~\ref{tab:replication} is empirical simulation on one dataset realization. A theoretical treatment would explain \emph{why} epsilon=0.1 (flip probability $\approx$47.5\% kept) still retains a mean gap of 0.012 rather than washing fully to zero, and would generalize past this specific sweep.
\item Even 30 replicates is a finite sample: the reported survival rates (e.g.\ 90.0\% at epsilon=4/5) carry their own sampling uncertainty (binomial SE $\approx$5.5\% at $n=30$, $p=0.9$) that is not itself propagated into the crossover-epsilon claim.
\item Single dataset (NIH ChestX-ray14, via Study A's frozen output), single label (Pneumothorax), single mechanism family (randomized response) --- other DP mechanisms for the same question (e.g.\ different local-DP constructions, or an aggregate-release design that avoids the discarded version's specific flaws) are not compared here.
\end{itemize}

\end{document}
